% p152_1; source PDF page 152; manually verified editable table.
\begin{tikzpicture}[x=1cm,y=-1cm,inner sep=0pt,outer sep=0pt]
\path[use as bounding box] (0,0) rectangle (16.000,8.370);
\draw[line width=.4pt] (0.000,0.000) rectangle (3.200,0.855);
\node[align=center,text width=2.940cm,font={\linespread{1}\fontsize{10}{13.50}\selectfont\bfseries}] at (1.600,0.427) {前提事实};
\draw[line width=.4pt] (3.200,0.000) rectangle (8.400,0.855);
\node[align=center,text width=4.940cm,font={\linespread{1}\fontsize{10}{13.50}\selectfont\bfseries}] at (5.800,0.427) {具体情形};
\draw[line width=.4pt] (8.400,0.000) rectangle (16.000,0.855);
\node[align=center,text width=7.340cm,font={\linespread{1}\fontsize{10}{13.50}\selectfont\bfseries}] at (12.200,0.427) {结论};
\draw[line width=.4pt] (0.000,0.855) rectangle (3.200,8.370);
\node[align=left,text width=2.940cm,font={\linespread{1}\fontsize{10}{13.50}\selectfont}] at (1.600,4.613) {甲准备拿着自己配的钥匙去盗窃丙家，乙知道后，送甲一把有用的钥匙，甲有了“双保险”的心理};
\draw[line width=.4pt] (3.200,0.855) rectangle (8.400,3.560);
\node[align=left,text width=4.940cm,font={\linespread{1}\fontsize{10}{13.50}\selectfont}] at (5.800,2.208) {备而不用。甲先用自己的钥匙，打开了门，窃得财物};
\draw[line width=.4pt] (8.400,0.855) rectangle (16.000,3.560);
\node[align=left,text width=7.340cm,font={\linespread{1}\fontsize{10}{13.50}\selectfont}] at (12.200,2.208) {乙的物理性帮助虽然没有发挥实际的帮助作用，但是起到心理性帮助效果，甲的心里有“双保险”的心理优势，故乙构成帮助犯既遂};
\draw[line width=.4pt] (3.200,3.560) rectangle (8.400,6.265);
\node[align=left,text width=4.940cm,font={\linespread{1}\fontsize{10}{13.50}\selectfont}] at (5.800,4.913) {弃而不用。甲先用乙的钥匙，一时打不开门，便用自己的钥匙，打开了门，窃得财物};
\draw[line width=.4pt] (8.400,3.560) rectangle (16.000,6.265);
\node[align=left,text width=7.340cm,font={\linespread{1}\fontsize{10}{13.50}\selectfont}] at (12.200,4.913) {乙的钥匙没有发挥实际的物理性帮助作用，也没有起到心理性的帮助效果，故乙不构成帮助犯既遂，而是未遂};
\draw[line width=.4pt] (3.200,6.265) rectangle (8.400,8.370);
\node[align=left,text width=4.940cm,font={\linespread{1}\fontsize{10}{13.50}\selectfont}] at (5.800,7.318) {有用。甲先用乙的钥匙，打开了门，窃得财物};
\draw[line width=.4pt] (8.400,6.265) rectangle (16.000,8.370);
\node[align=left,text width=7.340cm,font={\linespread{1}\fontsize{10}{13.50}\selectfont}] at (12.200,7.318) {乙的钥匙发挥了实际的物理性帮助，乙构成帮助犯既遂};
\end{tikzpicture}
